\section{Introducción: una «última entrevista» grabada antes de la adquisición}

Esta sección explica primero por qué este episodio merece una nota propia. Se grabó el 1 de diciembre de 2025. En el momento de la grabación, Meta aún no había anunciado la adquisición total de Manus; cuando el programa se publicó, ya se había convertido en la última entrevista de Manus. Este desfase temporal es importante porque le da a la entrevista una condición poco común: sin conocer el desenlace, el invitado expone de forma completa su trayectoria como emprendedor, su criterio de producto, su lógica de crecimiento, sus opiniones sobre la industria y sus preocupaciones personales.

El eje de este episodio no es la reconstrucción periodística de «una empresa que fue adquirida», sino cómo un equipo de producto de Agents pasó del desarrollo individual, el NLP, los grafos de conocimiento, Open IE, los GUI Agents y el crecimiento de Manus, hasta llegar a una filosofía de producto según la cual «la IA se parece más a la manufactura» y que «teme la complejidad». Puede enlazarse en una misma línea con Dokie en el EP130, la revisión general de Agents en el EP139 y la perspectiva de la salida a bolsa de empresas de modelos en el EP129: la capacidad de los modelos está entrando en el nivel del producto y de la organización de la empresa, y lo verdaderamente difícil es convertir esa capacidad en sistemas que se usen de forma sostenible, generen ingresos sostenibles y se mantengan simples de forma sostenible.

\begin{importantbox}{Tesis central del episodio}
La historia de Manus no es solamente la de un «Agent que se volvió un éxito viral», sino la de cómo empaquetar capacidades de IA en un sistema de productividad que los usuarios estén dispuestos a usar de forma continua y por el que estén dispuestos a pagar. Su mayor riesgo no es solo la competencia de las grandes empresas, sino perder su carácter distintivo y volverse complejo durante el proceso de crecimiento.
\end{importantbox}

\begin{knowledgebox}{Nota sobre la estrategia visual}
Este video es una entrevista con encuadre fijo, sin slides, pizarrón ni demostraciones de producto. El cuerpo de la nota no repite capturas de las personas; la portada sirve para identificar la fuente, y en el cuerpo se usan diagramas conceptuales, como la ruta de deriva hacia los Agents, la metáfora de la manufactura, el riesgo de la complejidad y el mapa del mercado, para explicar el contenido.
\end{knowledgebox}

\subsection{Resumen de la sección}

El EP128 es una entrevista que «cobra más sentido si se lee conociendo ya el desenlace». No solo cuenta cómo creció Manus, sino también cómo una empresa de Agents entiende la plataforma, el producto, la organización, la complejidad y la competencia futura.
